\section{Reference scales and the loss of charge exponent}\label{sec:vtm-bounds}

The reference side is chosen from the logarithmic charge
scale. Use $L,M,h$ from the introduction and the localization
power two. Thus
\begin{equation}\label{eq:transfer-scale-choice}
 e_K=e\sqrt h,\qquad v_K=-\log e_K,\qquad
 r_K=\min\{L^i:L^i\ge v_K^2,\ i\in\Z_{\ge0}\},
 \qquad R^*=\frac{hMr_K}{L}.
\end{equation}
Assume $v_K\ge8$, $e\le1/3$, $e_K\log h\le1$, and
\begin{equation}\label{eq:transfer-tail-condition}
 6N_*Mv_K^2 \exp(-v_K)
                 \left(v_K\exp(-v_K)+\tfrac12\right)\le\pi^2.
\end{equation}
These are finite charge conditions after $h,L,M$ are fixed.
They hold on a sufficiently small charge interval:
$v_K\to\infty$ as $e\to0$, and the polynomial times
exponential in \eqref{eq:transfer-tail-condition}
tends to zero.

Minimality gives $v_K^2\le r_K<Lv_K^2$.
Consequently $T^*=\log_L(Mr_K)-1$ is an integer at
least one, $R^*=hL^{T^*}$, and
$\ell_c=R^*/(2L)\to\infty$ as $e\to0$.
Moreover $1/L<\bar mh\le1$, so
\begin{equation}\label{eq:transfer-volume-error}
 \bar mR^*>\frac{Mr_K}{L^2}\ge r_K\ge v_K,
 \qquad \exp(-\bar mR^*)\le e\sqrt h.
\end{equation}
The same upper bound on $r_K$ gives
\begin{align*}
 e^2R^*\log(1/e)
 &<Mv_K^2 \exp(-2v_K)
                          (v_K+\tfrac12\log h)\\
 &\le Mv_K^2 \exp(-v_K)
                          (v_K\exp(-v_K)+\tfrac12).
\end{align*}
Thus \eqref{eq:transfer-tail-condition} implies
\eqref{eq:transfer-W3}.
The probe fits every sufficiently small charge and
sufficiently fine mesh, since $D_\beta^{\mathrm{buf}}$
is fixed and $\ell_c\to\infty$.

We now isolate the two reference estimates used in the
passage to arbitrary volumes. Their dependence on the
charge matters: a bounded partition defect still loses
a small power of $e$ through its photon contribution.
Let $C_{\mathrm{mat}}$ and $C_{\mathrm{ref}}$ be finite constants
independent of the charge, mesh and placement. Assume
that for $N\ge N_{\mathrm{ref}}(e)$ the mass assignment and reference
response satisfy
\begin{align}
 D^{\mathrm{mat}}&\le C_{\mathrm{mat}},\notag\\
 \Theta_4^{(0)}&\le C_{\mathrm{ref}} e^{\mathfrak e},\qquad
 (1+|c_2|)\Theta_2^{(0)}
                          \le C_{\mathrm{ref}} e^{\mathfrak e},
                                   \label{eq:transfer-terminal-input}
\end{align}
where
\begin{equation}\label{eq:transfer-charge-exponents}
 \mathfrak e=\tfrac12(1/4-9\varepsilon)-12\varepsilon,\qquad
 \sigma=\mathfrak e-\frac{2a_L}{N_*}>0,\qquad
 0<\varepsilon\le1/200.
\end{equation}
Here $N_{\mathrm{ref}}(e)$ is a fixed finite mesh threshold.
These are the reference hypotheses of \cref{thm:vtm-main}.
For the common tuned mass, the matter estimate is
\cite[Theorem~A, (5)]{TerminalPressure}, and the response
estimate is \cite[Theorem~A]{ReferenceResponse}.
Their joint charge interval and cutoff threshold are part
of the construction data assumed here.

On the reference torus take the free coefficients
$c_m=N_0^{(m)}(0)$ from
\eqref{eq:response-coefficient-relations}, using the
source associated with the same placed probe.
\cite[Proposition~5.2]{PositiveProbe}
and the support fit give a constant, enlarged into
$C_{\mathrm{ref}}$ if necessary, such that
\begin{equation}\label{eq:transfer-free-input}
 |c_2|+|c_4|^{1/2}\le C_{\mathrm{ref}},\qquad
 |c_4-3c_2^2-b(f)|
              \le C_{\mathrm{ref}}(e+\eta^{1/4}).
\end{equation}
Indeed $R^*\ge D_\beta+1$, and
\eqref{eq:transfer-volume-error} bounds the finite-volume
error in \cite[Theorem~B]{PositiveProbe}.
The coefficients are real and independent of placement.

\begin{proposition}\label{prop:transfer-UV-inputs}
Fix the probe and the scale data above, and assume
\eqref{eq:cells-fit}, \eqref{eq:transfer-terminal-input} and \eqref{eq:transfer-free-input}.
There is a finite
constant $C_U$, independent of $e,N,T$ and placement,
such that for all the stated reference charges and
$N\ge N_{\mathrm{ref}}(e)$,
\begin{equation}\label{eq:transfer-small-Wick}
 u_4\le C_Ue^\sigma,\qquad
 (1+|c_2|)u_2\le C_Ue^\sigma,\qquad
 |c_4-3c_2^2-b(f)|\le C_U(e^\sigma+\eta^{1/4}).
\end{equation}
On every cutoff torus with $T\ge T^*+1$, every finite
set $Y\subset(R^*/L)\Z^3$ of points incongruent
modulo $R\Z^3$ satisfies
\begin{equation}\label{eq:transfer-translated-Wick}
 \left|\E_{N,T}\prod_{y\in Y}
       W_m(F^\sharp(\tau_yf^\eta);v_\eta,c_m)\right|
                          \le u_m^{|Y|},\qquad m=2,4,
\end{equation}
Here $W_m$ is the polynomial in \eqref{eq:vtm-Wick}.
The fourth cumulant of the fixed smooth test obeys
\begin{equation}\label{eq:transfer-smooth-cumulant}
 |\kappa_4(F^\sharp(f))-e^4b(f)|
 \le(8C_U+3C_U^2)e^{4+\sigma}
                  +C_U(\eta^2+e^4\eta^{1/4}).
\end{equation}
The tests $f^\eta$ have one fixed compact support,
$f^\eta\to f$ in $L^2$, and $v_\eta\to\|f\|_2^2$.
Thus \eqref{eq:transfer-small-Wick} and
\eqref{eq:transfer-translated-Wick} give exactly the
cutoff hypotheses of Proposition~\ref{prop:UV-component-cumulant}.
\end{proposition}

\begin{proof}
Write
\[
 s_\beta=\|f\|_2+\|\beta\|_{H^4},\qquad
 C_\Gamma^{\mathrm{tr}}=(2\pi)^{a_L/N_*}
         \exp\left(\frac{K_L+C_{\mathrm{mat}}+1}
                                  {(L-1)^3N_*}\right).
\]
Since $R^*\ge1$, \eqref{eq:transfer-delta} gives
\begin{equation}\label{eq:transfer-Gamma-loss}
 \exp(\delta)\le C_\Gamma^{\mathrm{tr}} e^{-2a_L/N_*}.
\end{equation}
The inverse power here is a power of the charge.
The source-curvature identity and Jensen's inequality
for the disjoint averaging cubes give
$v_\eta^{1/2}\le\|f^\eta\|_2\le s_\beta$.
Use this and \eqref{eq:transfer-free-input} in
\eqref{eq:transfer-Wick-data}:
\begin{align*}
 T_4
 &\le C_{\mathrm{ref}} e^{\mathfrak e}
           +1200N_*^2s_\beta^4e^2+2C_{\mathrm{ref}}^2e^6\\
 &\le(C_{\mathrm{ref}}+1200N_*^2s_\beta^4+2C_{\mathrm{ref}}^2)e^{\mathfrak e},\\
 (1+|c_2|)T_2
 &\le C_{\mathrm{ref}} e^{\mathfrak e}
      +2e^6(12N_*s_\beta^2e^{-2}+C_{\mathrm{ref}})(1+C_{\mathrm{ref}})\\
 &\le[C_{\mathrm{ref}}+2(12N_*s_\beta^2+C_{\mathrm{ref}})(1+C_{\mathrm{ref}})]
                                                    e^{\mathfrak e}.
\end{align*}
Here $0<\mathfrak e<1/4$, from
\eqref{eq:transfer-charge-exponents}.
One possible constant is
\begin{equation}\label{eq:transfer-CU}
 \begin{split}
 C_U=\max\{&
 C_\Gamma^{\mathrm{tr}}(C_{\mathrm{ref}}+1200N_*^2s_\beta^4+2C_{\mathrm{ref}}^2),\\
 &C_\Gamma^{\mathrm{tr}}[C_{\mathrm{ref}}+2(12N_*s_\beta^2+C_{\mathrm{ref}})(1+C_{\mathrm{ref}})],
              \ C_{\mathrm{ref}},\ 60s_\beta^4\}.
 \end{split}
\end{equation}
Multiplication by \eqref{eq:transfer-Gamma-loss}
proves the first two bounds in
\eqref{eq:transfer-small-Wick}. The last follows from
\eqref{eq:transfer-free-input}, since $e\le e^\sigma$.

For a cutoff side $n=8hL^{T+N}$ with $T\ge T^*+1$,
one has $n=8L^{T-T^*}n_*$, a multiple of $p_{\rm cell}$ with
$n\ge Ln_*$. For $y\in(R^*/L)\Z^3$ use the even
cell index $c(y)$ of
\eqref{eq:cells-translation-observable}.
Incongruent points give distinct cells.
Proposition~\ref{prop:transfer-Wick} applies to their
product, and lattice translation invariance removes
the common placement. This proves
\eqref{eq:transfer-translated-Wick}.
For a fixed finite set of distinct points it applies
as soon as $R>2\max_{y\in Y}|y|_\infty$.

It remains to replace the discrete smooth test
$f^\eta$ by $f$ in the cumulant. We give the
quantitative multilinear estimate used for this purpose.
For a real two-cochain $g$, charge conjugation and
Gaussian domination give
\[
 \E\cosh(t\langle g,F^{\mathrm u}\rangle)
                          \le\exp(t^2\|g\|^2/2).
\]
Extract its nonnegative term of degree $2i$ and choose
$t^2=2i/\|g\|^2$. This gives
\[
 \E\langle g,F^{\mathrm u}\rangle^{2i}
    \le(2i)!\left(\frac{\|g\|^2}{2i}\right)^i \exp(i).
\]
For $i=2$ the bound is $(3/2)\exp(2)\|g\|^4$.
H\"older bounds the four-field moment by
\[
 \tfrac32\exp(2)\prod_a\|g_a\|.
\]
The variance bound bounds each of its three pair subtractions
by $\prod_a\|g_a\|$. Since $(3/2)\exp(2)+3<15$,
the symmetric multilinear forms $K_2,K_4$ obey
\begin{equation}\label{eq:transfer-multilinear-cumulants}
 |K_2(g_1,g_2)|\le\|g_1\|\|g_2\|,\qquad
 |K_4(g_1,g_2,g_3,g_4)|\le15\prod_a\|g_a\|.
\end{equation}
For $g=0$ the corresponding moment is zero.

Let $g$ and $g'$ be the periodizations of the normalized
covering cochains $f_{[\eta]}$ and $(f^\eta)_{[\eta]}$
from \cref{sec:reference-cells} onto the cutoff torus.
The common support cube has side less than $R$.
The averaging bound and \eqref{eq:cells-test-error} give
\[
 \|g\|,\|g'\|\le s_\beta,\qquad
 \|g-g'\|\le\eta^2\|\beta\|_{H^4}.
\]
The observables paired with these cochains are
$F^\sharp(f)$ and $F^\sharp(f^\eta)$.
Telescope the four slots of $K_4$ and the two slots
of $K_2$ in \eqref{eq:transfer-multilinear-cumulants}.
This proves
\begin{align}
 |\kappa_4(F^\sharp(f))-\kappa_4(F^\sharp(f^\eta))|
                              &\le60s_\beta^4\eta^2,\notag\\
 |\kappa_2(F^\sharp(f))-\kappa_2(F^\sharp(f^\eta))|
                              &\le2s_\beta^2\eta^2.
                                   \label{eq:transfer-test-replacement}
\end{align}
One-cell transfer and \eqref{eq:transfer-small-Wick}
bound the fourth cumulant of $F^\sharp(f^\eta)$ by
\[
 |\kappa_4(F^\sharp(f^\eta))-e^4(c_4-3c_2^2)|
                     \le(7C_U+3C_U^2)e^{4+\sigma}.
\]
Adding the free-coefficient error and the first bound
in \eqref{eq:transfer-test-replacement} proves
\eqref{eq:transfer-smooth-cumulant}.
Finally the compact support, $L^2$ convergence and
variance convergence are respectively
\eqref{eq:cells-smooth-test}, \eqref{eq:cells-test-error}
and \eqref{eq:cells-variance-limit}.
\end{proof}


